%% parts/p02/first-order-completeness-remainder.tex --- retained additions
\noindent\textit{Reader note.} The exact inherited chapter through source line 6048
replaces the earlier elementary-substructure and L\"owenheim--Skolem overlap.
The following separately attributed historical lines 314--371 retain only
the elementary-diagram and internal--external calibrations; they do not replace
or alter the inherited proofs.

\section{Compactness applications}
\label{sec:p02-first-order-compactness-applications}
%% Depth row D-P3-013: elementary-diagram upward LS, nonstandard models,
%% and the internal/external calibration of Skolem's paradox.

For the elementary-diagram construction and these model-theoretic
calibrations, see \cite{Hodges1993,Marker2002}.

\begin{hbdefinition}{def:p02-elementary-diagram}{Parameter expansion and elementary diagram}
For a \(\Lang\)-structure \(\mathcal M\), let
\(\Lang(\mathcal M)\) add a distinct constant \(c_a\) for each
\(a\in\dom\mathcal M\), and let \(\mathcal M_{\mathcal M}\) interpret
\(c_a\) as \(a\).  The \emph{elementary diagram}
\[
 \operatorname{ElDiag}(\mathcal M)
 \defeq
 \operatorname{Th}_{\Lang(\mathcal M)}(\mathcal M_{\mathcal M})
\]
is a theory---all parameter sentences true in that expansion---not a
structure.  Unlike the parameter-free theory
\(\operatorname{Th}_{\Lang}(\mathcal M)\), it records the truth of every
formula at each named tuple; that extra information is exactly what the
elementary-extension construction uses.
\index{elementary diagram}
\index{parameter expansion}
\end{hbdefinition}

\begin{hbtheorem}{thm:p02-upward-ls-elementary-extension}{proved=annex:p2.2.20}{Upward L\"owenheim--Skolem via elementary diagrams}
Assume the Axiom of Choice.  Let \(\mathcal M\) be an infinite
\(\Lang\)-structure and
\[
 \kappa\geq
 \max\bigl(\abs{\Lang},\abs{\dom\mathcal M},\aleph_0\bigr).
\]
Then there are a \(\Lang\)-structure \(\mathcal N\) of cardinality exactly
\(\kappa\) and an elementary embedding
\(j:\mathcal M\to\mathcal N\) whose range is proper.
\end{hbtheorem}

\begin{hbcorollary}{thm:p02-nonstandard-arithmetic-model}{proved=annex:p2.2.21}{A countable nonstandard elementary model of arithmetic}
Assume the Axiom of Choice.  Then the standard first-order structure
\(\mathcal N_{\mathrm{ar}}=(\N,0,S,+,\cdot)\) has a countable proper
elementary extension.  Consequently a countable model satisfies every
first-order sentence true in \(\mathcal N_{\mathrm{ar}}\) while containing
elements outside the embedded standard natural numbers.
\end{hbcorollary}

\begin{hbwarning}{def:p02-skolem-paradox}{Skolem's paradox is an internal--external distinction}
Assume the Axiom of Choice.  Conditionally, if a first-order set theory in a
countable signature has an infinite model, downward L\"owenheim--Skolem
supplies an externally countable elementary model.  Such a model can satisfy
that one of its sets is
\emph{uncountable}: internally there is no bijection belonging to the model,
while an external enumeration of the model need not be one of its elements.
There is no contradiction and no unconditional model-existence claim here;
the quantifiers range over different domains.  See
\cite{Skolem1923SetTheory,Hodges1993}.
\index{Skolem's paradox}
\end{hbwarning}
